\section{Method}
\label{sec:method}
% DRAFT (Claude): awaits the author's rewrite. Every setting below matches
% configs/pilot.yaml and src/markov_lean/{featurize,msm,judge}.py at the time
% of writing; recheck before submission.

\subsection{Proofs as trajectories}
% Data contract: docs/schema.md, table `steps`.
Each proof is a sequence of goal states $\state_0, \state_1, \dots, \state_n$,
where $\state_t$ is the pretty-printed goal before tactic $t$. Time is counted
in proof steps, so a lag $\lag$ and every timescale below are numbers of
tactics within one proof. Each proof is its own trajectory: no transition is
counted from the end of one proof to the start of the next.

We map each goal state to a vector with TF-IDF over Lean identifiers followed
by a truncated SVD to 128 dimensions and $\ell_2$ normalization. An
identifier such as \texttt{Nat.Prime.two\_le} contributes its full name, each
namespace prefix (\texttt{Nat}, \texttt{Nat.Prime}) and its parts split on
dots and underscores (\texttt{Prime}, \texttt{two}, \texttt{le}), so states
that share namespaces land close together. Single-character names and bare
numbers are dropped. \todo{if a sentence-embedding featurizer is added,
describe it and cite it here.}

\subsection{Microstates and the transition matrix}
% Table `assignments`, column `microstate`.
We cluster the state vectors into $n$ microstates with mini-batch $k$-means,
choosing $n \in \{100, 200, 400\}$ by the VAMP-2 score on held-out proofs.
At lag $\lag$ we count transitions $C_{ij}(\lag)$ from microstate $i$ at step
$t$ to $j$ at step $t+\lag$ with a sliding window within each proof, keep the
largest strongly connected set, and estimate a reversible maximum-likelihood
transition matrix $\T(\lag)$ with deeptime \citep{hoffmann2022deeptime}.

\subsection{Metastable areas}
% Table `assignments`, columns `macrostate`, `membership`.
PCCA+ \citep{deuflhard2005pcca,roblitz2013fuzzy} groups microstates into
$m$ metastable sets using the leading eigenvectors of $\T(\lag)$, and returns
a membership matrix $\chi \in [0,1]^{n \times m}$ whose rows sum to one. Each
step's macrostate is the set with the largest membership. We fix $m = 10$
before looking at any labels, and report agreement for
$m \in \{2, 4, \dots, 12\}$ alongside it.

\subsection{Validation}
The implied timescales
\begin{equation}
  t_i(\lag) = -\frac{\lag}{\ln |\lambda_i(\lag)|}
\end{equation}
are computed from the eigenvalues $\lambda_i$ of $\T(\lag)$. If the dynamics
are Markovian at lag $\lag$, the slow $t_i$ stop changing as $\lag$ grows; we
pick the smallest lag at which they do. The Chapman--Kolmogorov test then
checks that $\T(\lag)^k \approx \T(k\lag)$ on the macrostates
\citep{prinz2011markov}. Error bars come from a bootstrap over proofs with a
fixed seed. Because most traced proofs are short, we keep proofs with at
least 8 steps, so that a lag several steps long still leaves transitions to
count.

\subsection{LLM judges}
% Table `labels`.
No Lean source carries area labels, so two LLM judges, Claude Haiku 5.5 and
Claude Sonnet 5.5, independently assign a two-digit MSC 2020 code
\citep{msc2020} to each of 1{,}000 sampled proofs and to each of their steps.
We report Cohen's $\kappa$ between the judges and each judge's accuracy on
statements whose MSC code is known, and score macrostates against the
judges' consensus with normalized mutual information (NMI) and the adjusted
Rand index. The baseline partition labels every step of a proof with the
top-level Mathlib folder of its source file.
